% GENERATED, edit the French unit or the translation table
\mxcorrections{10}\label{C10-CORR}
\textbf{Teacher use.} The answers below concern the tasks not worked in the lesson. First have students state their method; give graded help before the solution. The durations are estimates to be calibrated in class.
\subsection*{Diagnostic and remediation}
\begin{enumerate}
\item $2\pi$. \textit{Action:} Link radius, diameter and perimeter.
\item $90^\circ$. \textit{Action:} Divide a circle into four sectors.
\item $\dfrac{\sqrt 2}{2}$. \textit{Action:} Rationalise and check the square.
\item $\pt{a;-b}$. \textit{Action:} Draw the two points.
\item $-\dfrac{\pi}{3}$. \textit{Action:} Put over a common denominator.
\item $\dfrac{3}{5}$. \textit{Action:} Identify the reference angle.
\end{enumerate}
\subsection*{Investigation activity}
Equilateral triangle then isosceles right triangle. Pythagoras gives respectively height $\dfrac{\sqrt 3}{2}$ and equal coordinates $\dfrac{\sqrt 2}{2}$. The signs are positive in the first quadrant.
\subsection*{Immediate applications}
\begin{enumerate}
\item $\dfrac{\pi}{4}$, $\dfrac{2\pi}{3}$; length $\dfrac{3\pi}{2}$ m.
\item $\dfrac{7\pi}{2}$ has the image of $\dfrac{3\pi}{2}$; $-\dfrac{5\pi}{3}$ that of $\dfrac{\pi}{3}$.
\item $\sin ^2x=\dfrac{16}{25}$; the sine is negative, so $-\dfrac{4}{5}$.
\item $\pt{-\dfrac{1}{2};-\dfrac{\sqrt 3}{2}}$ and $\pt{\dfrac{\sqrt 2}{2};-\dfrac{\sqrt 2}{2}}$.
\item $h=\dfrac{\sqrt 3}{2}$; new chord $\dfrac{1}{\sqrt {2+\sqrt3}}=\sqrt {2-\sqrt3}$.
\end{enumerate}
\subsection*{Exercises: method and results}
\paragraph{Exercise 10.1 -- Converting and measuring.} $\dfrac{5\pi}{6}$ and $-\dfrac{\pi}{2}$; arc $\dfrac{4\pi}{3}$ cm. The negative sign indicates an orientation; a geometric length is positive.
\paragraph{Exercise 10.2 -- Reducing angles.} $\dfrac{5\pi}{6}$, $\dfrac{\pi}{4}$, $\pi$. Remove or add multiples of $2\pi$.
\paragraph{Exercise 10.3 -- Exact values.} $\pt{-\dfrac{1}{2};\dfrac{\sqrt 3}{2}}$; $\pt{\dfrac{\sqrt 3}{2};-\dfrac{1}{2}}$; $\pt{-\dfrac{\sqrt 2}{2};-\dfrac{\sqrt 2}{2}}$, in the form (cosine; sine).
\paragraph{Exercise 10.4 -- Recovering the sign.} $\cos ^2x=1-\dfrac{25}{169}=\dfrac{144}{169}$; cosine negative, so $-\dfrac{12}{13}$.
\paragraph{Exercise 10.5 -- A wheel.} Advance $r\theta =\num{0,3}(\dfrac{5\pi}{3})=\dfrac{\pi}{2}$ m. This corresponds to the unrolling of the arc; the chord joins two positions on the circle and does not have this length. A material point on the wheel follows a different path relative to the ground.
\paragraph{Exercise 10.6 -- Chord and scalar product.} $A\pt{1;0}$, $B\pt{-\dfrac{1}{2};\dfrac{\sqrt 3}{2}}$; $AB^2=\dfrac{9}{4}+\dfrac{3}{4}=3$. Al-Kashi: $1+1-2\cos (\dfrac{2\pi}{3})=3$. So $AB=\sqrt3$.
\paragraph{Exercise 10.7 -- Reading an equation on the circle.} The vertical line with x-coordinate $\dfrac{1}{2}$ meets the circle at the angles $\dfrac{\pi}{3}$ and $\dfrac{5\pi}{3}$. The two points have opposite sines.
\paragraph{Exercise 10.8 -- Delayed review.} $\dfrac{25\pi}{4}=6\pi +\dfrac{\pi}{4}$. Three complete turns do not change the image, so both values are $\dfrac{\sqrt 2}{2}$.
\paragraph{CQFD, exercise 32 p.~220, questions a and b.} a: $\cos(\dfrac{6\pi}5)=\cos(\pi+\dfrac\pi5)=-\dfrac{1+\sqrt5}{4}$. b: $\cos(-\dfrac\pi5)=\cos(\dfrac\pi5)=\dfrac{1+\sqrt5}{4}$, by evenness of the cosine.
\paragraph{CQFD, exercise 39 p.~220, questions 1, 2, 3 and 5.} 1: $\dfrac\pi6+\dfrac\pi3=\dfrac\pi2$. 2: yes, $\cos(\dfrac\pi3)=\dfrac12=\num{0,5}$ is a decimal. 3: the principal measure in $\intof{-\pi}{\pi}$ is $-\dfrac{5\pi}6$. 5: $\sin(\dfrac{4\pi}3)=-\dfrac{\sqrt3}2$.
\subsection*{MCQ: answers and reasons}
\textbf{MCQ 1.}
\begin{enumerate}
\item B. Half turn.
\item A. Equilateral triangle.
\item D. Lowest point of the circle.
\item C. Equation of the unit circle.
\item B. Complete turns.
\end{enumerate}
\textbf{MCQ 2.}
\begin{enumerate}
\item C. Vertical symmetry.
\item A. Isosceles right triangle.
\item D. Three quarters of a turn.
\item B. $r\theta=6$.
\item C. Negative x-coordinate inside the quadrant.
\end{enumerate}
\subsection*{True or false}
\begin{enumerate}
\item False: at 0, the two sides would be 1 and $-1$. The cosine is even.
\item True: a complete turn keeps the image.
\item False: $\dfrac{\pi}{6}$ and $\dfrac{5\pi}{6}$ have the same sine but opposite cosines.
\end{enumerate}
\paragraph{Differentiated teaching.} For a technical block, offer only the first step (domain, factor or reference sample space), then withdraw the help. For an independent student, ask for a second method and the treatment of a limiting case. Orally, allow an initial reformulation in the language the student has mastered, then require a precise mathematical sentence in French.
